\documentclass[11pt]{article}
\usepackage{times}
\usepackage{fullpage}
\usepackage{url}
\usepackage{hyperref}
\usepackage{fancyhdr}
\usepackage{graphicx}
\usepackage{enumitem}
\usepackage{subcaption}
\usepackage{amsmath, amsfonts, amsthm, fouriernc}
\usepackage{IEEEtrantools}
\usepackage{parskip}

\pagestyle{fancy}
\addtolength{\headheight}{2em}
\addtolength{\headsep}{1.5em}
\lhead{\large{ME 2801}}

\usepackage{sectsty}
\sectionfont{\fontsize{12}{8}\selectfont}

\begin{document}

\begin{center}
  \Large{\bf{Assignment: Transfer Functions and Modeling}}
\end{center}

\textbf{Instructions.}  Show your work for full credit.  All transfer functions
are evaluated at zero initial conditions.  Express answers as ratios of
polynomials in $s$ in descending order.  MATLAB may be used to verify results
but hand derivation is required.

\bigskip\hrule\bigskip

% ===========================================================================
\section*{Part 1: Transfer Functions from ODEs}
% ===========================================================================

For each problem: (a)~identify the input and output variable; (b)~take the
Laplace transform of the ODE assuming zero initial conditions; (c)~solve for
the transfer function $G(s) = C(s)/R(s)$.

\begin{enumerate}[label=\textbf{\arabic*.}]

\item First-order, simple.
\[
  3\dot{c}(t) + c(t) = r(t)
\]

\item First-order, input derivative.
\[
  2\dot{c}(t) + 4c(t) = 5\dot{r}(t)
\]

\item Second-order, no input derivative.
\[
  \ddot{c}(t) + 6\dot{c}(t) + 8c(t) = r(t)
\]

\item Second-order, input appears on both sides.
\[
  \ddot{c}(t) + 3\dot{c}(t) + 2c(t) = \dot{r}(t) + 4r(t)
\]

\item Third-order.  Identify the poles.  Are any of them real?  What does that
  tell you about the step response?
\[
  \dddot{c} + 6\ddot{c} + 11\dot{c} + 6c = r(t)
\]

\item \textbf{Reverse direction.}  A system has transfer function
\[
  G(s) = \frac{s + 3}{s^2 + 4s + 3}
\]
Write the corresponding differential equation relating output $c(t)$ to input
$r(t)$.

\item Given $\ddot{c} + 5\dot{c} + 6c = 2\dot{r} + r$:
(a)~Find $G(s)$.
(b)~Find the poles and zeros of $G(s)$.
(c)~Use MATLAB to create the \texttt{tf} object and verify with \texttt{pole(G)}
and \texttt{zero(G)}.

\end{enumerate}

\bigskip\hrule\bigskip

% ===========================================================================
\section*{Part 2: Mechanical System Transfer Functions}
% ===========================================================================

For each system: (a)~draw a free-body diagram; (b)~write the equation of motion
using Newton's second law; (c)~take the Laplace transform; (d)~find the transfer
function with the input and output as stated.

\begin{enumerate}[label=\textbf{\arabic*.}, resume]

\item A mass $M = 2\,\text{kg}$ is connected to a wall by a spring $K = 8\,\text{N/m}$
  (no damping).  Input: applied force $f(t)$; output: displacement $x(t)$.

  Find $G(s) = X(s)/F(s)$.  What are the poles?  What type of motion does the
  step response produce?

\item A mass $M$ is connected to a wall by a spring $K$ and a viscous damper $b$
  in parallel.  Input: applied force $f(t)$; output: displacement $x(t)$.

  (a)~Find $G(s) = X(s)/F(s)$ in terms of $M$, $b$, $K$.

  (b)~Let $M = 1\,\text{kg}$, $b = 4\,\text{N{\cdot}s/m}$, $K = 3\,\text{N/m}$.
  Find the numerical transfer function, then use MATLAB to find the poles.  Is
  the system overdamped, underdamped, or critically damped?

\item A rotating disk has moment of inertia $J = 0.5\,\text{kg{\cdot}m}^2$
  and is connected to a fixed support through a rotational damper
  $D = 2\,\text{N{\cdot}m{\cdot}s/rad}$.  Input: applied torque $T(t)$; output:
  angular velocity $\omega(t) = \dot{\theta}(t)$.

  (a)~Write the equation of motion in terms of $\omega(t)$.

  (b)~Find $G(s) = \Omega(s)/T(s)$.

  (c)~Use the final value theorem to find the steady-state angular velocity for
  $T(t) = T_0\,u(t)$.

\item A rotating shaft has inertia $J$, torsional spring $K$, and rotational
  damping $D$.  Input: applied torque $T(t)$; output: angular displacement
  $\theta(t)$.

  Find $G(s) = \Theta(s)/T(s)$ in terms of $J$, $D$, $K$.

\end{enumerate}

\bigskip\hrule\bigskip

% ===========================================================================
\section*{Part 3: Using the Transfer Function}
% ===========================================================================

\begin{enumerate}[label=\textbf{\arabic*.}, resume]

\item A first-order system has transfer function $G(s) = \dfrac{5}{s+5}$.

  (a)~For a unit step input $R(s) = 1/s$, write $C(s) = G(s)R(s)$ and find
  $c(t)$ via partial fractions.

  (b)~Use the final value theorem to find the steady-state value without
  inverting the transform.

\item A second-order system has transfer function
\[
  G(s) = \frac{4}{s^2 + 2s + 4}
\]

  (a)~For a unit step input, write $C(s)$.

  (b)~Use MATLAB's \texttt{step(G)} to plot the step response.  Label the plot
  with the natural frequency $\omega_n$, damping ratio $\zeta$, and time
  constant $\tau = 1/(\zeta\omega_n)$.
  (Hint: match the denominator to $s^2 + 2\zeta\omega_n s + \omega_n^2$.)

\end{enumerate}

\end{document}
